% =============================================================================
% Results sec 5.2 -- Statistical heterogeneity alone gives limited separation
% Compressed version. All numerics verified against VERIFICATION_SHEET.txt.
% Target ~500 words.
% =============================================================================

\section{Statistical heterogeneity alone gives limited optimiser separation}
\label{sec:statistical-heterogeneity}

Under an IID partition (Flower, $n=10$ paired seeds), FedAvg and FedProx
differ by $\Delta = -0.007$ macro-F1 (FA $0.585$, FP $0.578$, FN $0.577$) ---
within seed noise. We treat this as a negative control: statistical
heterogeneity is a necessary condition for any FedAvg--FedProx gap on this
benchmark.

On the engineered balanced-paired seven-client partition we report three
estimates ($n=10$ paired seeds each): Flower paired $\Delta = +0.007$ (7/10
paired wins, full \texttt{git\_commit} and \texttt{client\_update\_norms}
provenance), Flower $\mu$-sweep peak at $\mu = 0.01$ $\Delta = +0.013$ (full
provenance, $\mu$ characterisation in \S\,\ref{sec:mu-sensitivity}), and
pure-PyTorch reference $\Delta = +0.027$ (9/10 wins, no \texttt{git\_commit},
no companion update-norm logs). All three measure the same partition through
slightly different code paths; they agree on sign and bracket the magnitude
at $0.01$--$0.03$. The Flower paired $\Delta = +0.007$ is the central
headline; the pure-PyTorch number is cross-runtime corroboration.

The pure-PyTorch $+0.027$ is concentrated in two classes: melanoma
contributes $+0.114$ per-class F1 (61\,\% share of the macro-F1 advantage)
and actinic keratoses $+0.067$ (36\,\%); the other five classes contribute
the residual 3\,\% (all $|\Delta| \leq 0.013$), and mel-nevi is invariant at
$+0.002$. Figure~\ref{fig:engineered-per-class} renders the per-class
validation trajectories under the engineered partition ($n=10$ paired
seeds, $\pm 1$ SEM bands): the melanoma and actinic panels show non-
overlapping FedAvg vs FedProx bands by mid-training, while the other five
classes (including mel-nevi) overlap throughout 150 rounds. The gain
reflects rebalancing of two minority classes, not a uniform optimiser
improvement.

Three other partitions replicate the small-magnitude finding: Dirichlet
$\alpha = 0.1$ $\Delta = +0.008$ ($n=10$), specialist 1-of-7 $\Delta = +0.008$
($n=10$), multi-seed node-pinned L4 $\Delta = +0.005$ ($n=3$; per-seed
$\{+0.013, +0.019, -0.016\}$). The single-seed heterogeneity-ladder pilot
at L4 gives $\Delta = -0.030$, but its $-0.030$ deficit is driven entirely
by vascular F1 dropping by $-0.220$ (FA $0.702$ $\rightarrow$ FP $0.482$);
all six other classes are within $|\Delta| \leq 0.034$. The aggregate is one
rare-class outcome on one seed, and the $n=3$ node-pinned promotion
($\Delta = +0.005$) averages it out. All Dirichlet aggregates above use
regenerated raw-JSON summaries; the previously-committed analysis CSVs were
stale and showed a spurious seed-8675309 catastrophic FedProx failure that
the regenerated data do not support.

On the engineered partition, FedProx reduces test-macro-F1 SD from $0.025$
to $0.014$ ($-45\,\%$), reaches validation macro-F1 $= 0.40$ in $9.8$ vs
$12.7$ rounds ($1.3\times$ speed-up), and shows a paired mid-training peak
$\Delta = +0.094$ at round 50 that closes to $+0.007$ by round 150. These
trajectory characteristics are documented for the engineered partition only.

Statistical heterogeneity alone produces small, runtime-sensitive separation
(Table~\ref{tab:stat-het-headline}, Figure~\ref{fig:stat-het-forest}). The
next question is whether the proximal coefficient $\mu$ itself explains when
FedProx helps.


% =============================================================================
% Table 5.2 placeholder
% =============================================================================

\begin{table}[h]
\centering
\caption{Statistical heterogeneity and runtime sensitivity. Six cells:
IID, the engineered partition through two runtimes plus the $\mu$-peak,
Dirichlet $\alpha=0.1$, specialist 1-of-7, and the multi-seed node-pinned
L4 control. Dirichlet, specialist, and IID values from regenerated raw-JSON
aggregates.}
\label{tab:stat-het-headline}
\resizebox{\textwidth}{!}{%
\begin{tabular}{lccccccl}
\hline
Setting              & Runtime       & Clients & $n$ & FedAvg              & FedProx             & $\Delta$ & Note \\
\hline
IID                  & Flower        & 7       & 10  & $0.585 \pm 0.020$   & $0.578 \pm 0.020$   & $-0.007$ & negative control \\
Engineered (pure)    & pure-PyTorch  & 7       & 10  & $0.481 \pm 0.025$   & $0.508 \pm 0.014$   & $+0.027$ & 9/10 wins; partial provenance \\
Engineered (Flower)  & Flower        & 7       & 10  & $0.496$             & $0.503$             & $+0.007$ & 7/10 wins; \textbf{central} \\
Dirichlet $\alpha=0.1$ & Flower      & 7       & 10  & $0.480 \pm 0.036$   & $0.488 \pm 0.034$   & $+0.008$ & regenerated raw-JSON \\
Specialist 1-of-7    & Flower        & 7       & 10  & $0.437 \pm 0.024$   & $0.445 \pm 0.017$   & $+0.008$ & --- \\
Node-pinned L4       & Flower        & 2       & 3   & $0.492 \pm 0.019$   & $0.497 \pm 0.022$   & $+0.005$ & multi-seed L4 control \\
\hline
\end{tabular}%
}
\end{table}


% =============================================================================
% Figure 5.2 placeholder
% =============================================================================

\begin{figure}[h]
\centering
\includegraphics[width=0.85\textwidth]{F_statistical_heterogeneity_forest.pdf}
\caption{Forest plot of $\Delta(\textrm{FedProx} - \textrm{FedAvg})$ macro-F1
across the six settings in Table~\ref{tab:stat-het-headline}. Whiskers show
$\pm 1$ sample SD across paired seeds; dashed line at $\Delta = 0$ marks no
separation. The IID negative control sits at the left; the five partitioned
settings cluster between $+0.005$ and $+0.027$.}
\label{fig:stat-het-forest}
\end{figure}


% =============================================================================
% Figure 5.2(b) -- engineered partition per-class trajectories
% =============================================================================

\begin{figure}[h]
\centering
\includegraphics[width=\textwidth]{F_engineered_per_class.pdf}
\caption{Per-class validation F1 trajectories on the engineered partition,
$n = 10$ paired seeds (FedAvg vs FedProx, mean $\pm 1$ SEM bands).
Two-by-four grid (seven classes plus legend). Melanoma and actinic-keratoses
panels show non-overlapping FedAvg/FedProx bands by mid-training --- the
two classes that carry the pure-PyTorch $+0.027$ macro-F1 advantage
($+0.114$ and $+0.067$ per-class F1, $61\,\%$ and $36\,\%$ of the share
respectively). The other five classes show overlapping bands throughout
150 rounds, including the dominant mel-nevi panel which is invariant near
its prevalence saturation.}
\label{fig:engineered-per-class}
\end{figure}
